\section{Mezclas químicas}



\section{Propiedades de las mezclas químicas}



\section{Concentraciones}

\definition{Concentración}{Es la proporción o relación que hay entre la cantidad de soluto y la cantidad de disolución.}

\definition{Soluto}{Sustancia que se disuelve en el solvente para formar una solución o disolución}

La concentración es una característica cuantitativa, o sea, que se puede medir numéricamente. Algunas de estas formas cuantitativas de medir la concentración son:

\begin{itemize}
    \item Molaridad
    \item Molaridad
    \item Porcentajes de soluto
    \item Partes por millón (ppm)
\end{itemize}

\subsection{Molaridad}

Aquellas en las que la masa del soluto está expresada en moles por litro de solución.

\definition{Mol}{1 mol de cualquier sustancia contiene $6.022 \times 10^{23}$\footnote{Este número es el famoso Número de Avogadro} entidades elementales de esa sustancia.}

Podemos comparar la unidad de medida del mol con otras formas de medida como la docena. Así como una docena de, por ejemplo huevos, significa que tenemos 12 huevos, 1 mol de huevos significaría que tendremos $6.022 \times 10^{23}$ huevos. Claramente esa cantidad de huevos es ligeramente elevada, y es por ello que no usamos el mol para medir huevos, sin embargo, cuando hablamos de moléculas si es adecuado usar el mol, pues la cantidad de átomos de un elemento dado en un simple gramo de algo es enorme.

\begin{align*}
\text{Molaridad} &= \frac{\text{moles del soluto}}{\text{litros de solución}}
\\
\\
\text{moles del soluto} &= \frac{\text{g de soluto}}{\text{peso molecular del soluto}}
\end{align*}

\subsection{Normalidad}

Aquellas en las que en 1 litro de agua hay disuelto un equivalente del soluto. Es decir, el peso molecular de la sustancia dividido por el número de electrones que intercambia en la reacción.

\begin{align*}
    N &= \frac{\text{Equivalentes de soluto}}{\text{Litros de solución}}
    \\
    \\
    \text{Equivalente de soluto} &= \frac{\text{g de soluto}}{\text{peso equivalente de soluto}}
    \\
    \\
    \text{Peso equivalente de soluto} &= \frac{\text{peso molecular de soluto}}{\text{parámetro de carga de soluto}}
\end{align*}

\subsection{Porcentajes de soluto}

Aquellas cuya medida es la cantidad de soluto por cada 100 unidades de solución (no de solvente).

\begin{align*}
    \%\frac{p}{v} &= \frac{\text{gramos de soluto}}{\text{volumen de solución}} \times 100
    \\
    \\
    \%\frac{p}{p} &= \frac{\text{gramos de soluto}}{\text{gramos de solución}} \times 100
    \\
    \\
    \%\frac{v}{v} &= \frac{\text{volumen de soluto}}{\text{volumen de solución}} \times 100
\end{align*}\footnote{p = peso; v = volumen}

\subsection{Partes por millón (ppm)}

Se refiere a la cantidad de unidades del soluto que hay por cada millón de unidades del solvente. Se abrevia como "ppm”.

Es un concepto análogo (igual) al de porcentaje. Mientras que en el porcentaje hablamos de cuantas unidades de soluto tienes por cada 100 unidades de solvente, en las ppm hablamos de cuantas partes de soluto tienes por cada millón de partes de solvente. De hecho, se puede definir la siguiente equivalencia:

\begin{equation*}
    10,000ppm = 1\%
\end{equation*}

En el caso de disoluciones acuosas, una parte por millón (1 ppm) equivale a un miligramo de soluto por cada litro de disolución. O si usamos unidades más pequeñas, un microgramo de soluto por mililitro de disolución:

\begin{align*}
    1ppm &= \frac{1mg}{1L}
    \\
    \\
    1ppm &= \frac{1\mu g}{1ml}
\end{align*}
